% -----------------------------------------------------------------------------------
\section{Register of stipulations}\label{app:stipulations}

The results are theorems about the definitions displayed in the text. The following choices are made
there, and each could be made differently; where a different choice is known to change a result, the
section that discusses it is given.

\begin{enumerate}[label=\arabic*.,leftmargin=2em,itemsep=.3em]
  \item Policies are deterministic unless a mixed variant is stated (Sections~\ref{sec:robust},
        \ref{sec:mixed}).
  \item Every admissible policy induces a well-defined probability law (Definition~\ref{def:welldefined});
        self-referential problems that lack one are excluded (Remark~\ref{rem:skalse}).
  \item No information set occurs twice on a path (Definition~\ref{def:problem}).
  \item Structural environment plans, including responses on branches not taken, are part of the problem
        (Definition~\ref{def:structural}); how outcomes on alternative branches are coupled is fixed by
        them.
  \item Causal deviations freeze the baseline distribution over plans, and a local deviation keeps the
        policy's other prescriptions (definition of $Q_\pi$).
  \item Dominance quantifies over structural states, including those of probability zero under the
        evaluated policy (Section~\ref{sec:criteria}); with supported states only, \cri{2} becomes
        unsatisfiable in Death in Damascus.
  \item Positive reach is measured under the evaluated policy.
  \item Predictor error is given by a kernel; the Newcomb calculations use symmetric error.
  \item A predictor targets specified policy coordinates; \enquote{the same computation} is represented
        by a stipulated dependence (Sections~\ref{sec:twin}, \ref{sec:bomb}).
  \item Commitment is irrevocable, precedes the environment's response, offers the full policy set, and
        uses a specified coupling and structural state set (Definition~\ref{def:commit},
        Section~\ref{sec:couplings}).
  \item Commitment extensions are tagged and well founded, and the class is closed under them.
  \item \cri{5} is the totality already built into the definition of a theory.
  \item The evidential criterion uses the reference-policy law and the null-event convention displayed in
        Section~\ref{sec:evidential}; other evidential semantics change the map (Section~\ref{sec:evsem}).
  \item Mixed variants specify sampling, predictor access, information, counterfactual coupling, and
        continuity (Sections~\ref{sec:mixedexist}, \ref{sec:mixed}).
  \item Cross-problem rules output a reflexive relation whose maximal elements are required to be nonempty
        (Definition~\ref{def:rule}); neutrality acts by bijections between label sets, and invariances act
        on the relation.
  \item Ordering concerns the output relation at each input, not rationalisability of each choice set.
  \item Menu restriction is principal unless stated otherwise; principal restriction is forced by
        law-preserving removal, and column and fixed-reference restriction require additional data
        (Section~\ref{sec:menus}).
  \item Non-triviality is existential unless a universe is specified (Theorem~\ref{thm:five}).
  \item On restricted domains, an invariance applies when both a matrix and its transform are admissible
        (Section~\ref{sec:translations}).
  \item Equivalence of problems and parameterisation of rules are made precise before taxonomy claims are
        stated (Section~\ref{sec:taxonomy}).
\end{enumerate}
